% Supplementary Note 6: the pooled analysis of second works. Numbers through macros only.
\suppnote{The pooled analysis of second works}

This note reports in full the exploratory analysis that pooled Studies 2, 4 and 5, which the main text summarizes in its Results section on the residue at slight edits. It gives the complete estimates, the rival accounts and the open explanations of the peak at Slight. It also reports two registered results from the separate studies that bear on how the penalty changes across levels: Study~2's test of dilution and Study~4's tests at each level.

\paragraph*{Design and measures.}
Studies 2, 4 and 5, with \STwoN, \SFourN{} and \SFiveN{} participants, asked people to rate second works that Alex was said to have made with his design and sketching tools from original images, all of which had been generated with an AI image generator. A label, varied between participants, said that Alex had made each original either with design and sketching tools (the Human label) or by prompting an AI system (the AI label). Each original came with three second works, and every participant saw all three: None, the original shown again unchanged; Slight, a light edit; and Dramatic, a heavy one. Participants rated each second work's transformativeness (how much it adds new expression, meaning or message to the original), its essence change (how much the essence of the original changed) and the creativity involved in its production, each on a 1-to-5 scale. The label's penalty is the mean rating under the Human label minus the mean under the AI label, in points on that scale; a positive penalty means that the AI label lowered the rating. We call what the label carried into ratings of second works the residue. Before any pooling, each study's mean ratings under the two labels at each level, and the penalty between them, are plotted in main-text Fig.~2b for creativity and in Supplementary Fig.~\ref{figS:otherratings} for the other two measures. Study~5's tables give 90\% intervals; the 95\% intervals plotted for it are computed from the same estimates and standard errors (a standard error is the expected size of an estimate's error).

\begin{figure}[tbp]
\centering
\makebox[\textwidth][c]{\includegraphics[width=183mm]{figS_other_ratings}}
\caption{\textbf{Transformativeness and essence change of the second work under the two labels.} Marks and columns are as in main-text Fig.~2b: the two label groups' mean ratings, as predicted by each study's statistical model, and the gap between them (the Human-label mean minus the AI-label mean) with its 95\% confidence interval. The label described only the original: both groups read that Alex made the second works with his design and sketching tools. \textbf{a},~Transformativeness, how much the second work adds new expression, meaning or message to the original (a lay rating, not a fair-use judgment). \textbf{b},~Essence change, how much the original's essence changed (defined in Studies 4 and 5 as its core sense and feeling). None is the original shown again unchanged. As in main-text Fig.~2b, each person saw every level but only one label, so the rise across levels and the gap cannot be compared in size. The gaps are exploratory in Study~2 and registered secondary tests in Studies 4 and 5, each judged at .05 without correction; they average over the hours stated for the original (Study~4) or over stated editing hours (Study~5). Supplementary Table~\ref{tab:supp-labeleffects} lists the registered tests.}
\label{figS:otherratings}
\end{figure}

\paragraph*{Status and model.}
This analysis is exploratory. It was specified on 9 October 2026, after every registered result was known; its tests are not corrected for multiple comparisons, and it changes no registered verdict. A $p$-value is the chance of a difference at least this large if there were none, and a result is called significant when $p$ is below .05. Where this note cites a study's own registered tests, each keeps the status its registration gave it: `registered confirmatory' marks a test that the registration names as confirmatory, and `registered secondary' marks one it named outside that set and judged at .05 without correction. A pooled model fitted Studies 2, 4 and 5 together, one measure at a time. It was a linear mixed model, a regression that gives each participant, each original and each second work its own baseline, and it estimated the label's penalty at each level with each study weighted equally, whatever its size. It left out Study~4's stated hours for the original and Study~5's stated editing hours, so each study's penalty averages over them. It also treats each study's participants as different people, so a person who took part in Study~5 and in Study~2 or~4 counts as two (Section~\ref{sec:suppB-study5}). Its fits converged, meaning that the fitting procedure settled on a solution, but they are singular: one variance, the spread between originals once each second work had its own baseline, was estimated at zero.

\paragraph*{The pooled penalties.}
On all three measures the pooled penalty was largest at Slight (Supplementary Fig.~\ref{figS:pooled}), though not in every study. On creativity it was \XSCreatPenaltyNoneEst, \XSCreatPenaltySlightEst{} and \XSCreatPenaltyDramaticEst{} points at None, Slight and Dramatic; on transformativeness \XSTransPenaltyNoneEst, \XSTransPenaltySlightEst{} and \XSTransPenaltyDramaticEst; and on essence change \XSEssPenaltyNoneEst, \XSEssPenaltySlightEst{} and \XSEssPenaltyDramaticEst. Only the Slight penalties differed from zero. Their 95\% confidence intervals (CIs), the ranges of values the data do not rule out at the 95\% level, ran from \XSCreatPenaltySlightLo{} to \XSCreatPenaltySlightHi{} on creativity, \XSTransPenaltySlightLo{} to \XSTransPenaltySlightHi{} on transformativeness and \XSEssPenaltySlightLo{} to \XSEssPenaltySlightHi{} on essence change. Ratings at None sat near the floor of the scale, where a penalty has little room to show, so a penalty near zero there says little. Studies 2, 4 and 5 registered a bound of $\pm$\SourceBound{} points for the penalty on second works, half of Study~1's creativity penalty on originals. An equivalence test asks whether a penalty's 90\% CI lies wholly inside that bound; a pass means that the penalty is smaller than the bound, not that it is zero. Read this way, the pooled penalty lay inside the bound at every level on creativity and transformativeness, and at None and Dramatic on essence change; at Slight the essence-change interval (\XSEssPenaltySlightLoNinety{} to \XSEssPenaltySlightHiNinety) reached past it. For reference, the creativity penalty on originals was \XSOriginalsSOneEst{} points in Study~1 (95\% CI \XSOriginalsSOneLo{} to \XSOriginalsSOneHi; registered confirmatory) and \XSOriginalsSThreeEst{} in Study~3 averaged over stated time (\XSOriginalsSThreeLo{} to \XSOriginalsSThreeHi; an exploratory recomputation from Study~3's registered model). Each pooled creativity penalty on second works was smaller than both, a reading across studies, not a test.

\paragraph*{Three rival accounts.}
Three accounts predict different shapes across the levels. Under taint, the AI origin passes to anything built from it, so the penalty is the same at every level. Under dilution, the penalty is largest at None and shrinks as the edit grows. Under severance, the label does not reach second works at all, so there is no penalty at any level. Two quantities summarize the shape. The slope is the change in penalty per step from None to Dramatic (half of the Dramatic penalty minus the None penalty), which dilution predicts to be negative. The curvature is the Slight penalty minus the average of the None and Dramatic penalties, which is zero for any straight line. The pooled slope did not differ from zero on any measure: \XSCreatSlopeEst{} points on creativity (95\% CI \XSCreatSlopeLo{} to \XSCreatSlopeHi), \XSTransSlopeEst{} on transformativeness (\XSTransSlopeLo{} to \XSTransSlopeHi) and \XSEssSlopeEst{} on essence change (\XSEssSlopeLo{} to \XSEssSlopeHi). The fall from None that dilution predicts is therefore not shown, though the floor at None weakens this test. Instead, on all three measures the penalty at Slight lay above the average of the penalties at None and Dramatic: by \XSCreatCurvatureEst{} points on creativity (\XSCreatCurvatureLo{} to \XSCreatCurvatureHi), \XSTransCurvatureEst{} on transformativeness (\XSTransCurvatureLo{} to \XSTransCurvatureHi) and \XSEssCurvatureEst{} on essence change (\XSEssCurvatureLo{} to \XSEssCurvatureHi).

\paragraph*{Registered tests in the separate studies.}
Two registered results from the separate studies bear on how the penalty changes across levels. Study~2 tested dilution directly as its third hypothesis (H3, registered confirmatory): the penalty at None minus the penalty at Dramatic, which dilution predicts to be positive. The test was judged after Holm's correction across Study~2's confirmatory tests, a correction\cite{holm1979simple} that raises each $p$-value in a set of tests so that the chance of calling any result in the set significant by mistake stays at .05. It was not significant on any measure, and every estimate ran against dilution: \STwoTransPenaltyChangeEst{} points on transformativeness (95\% CI \STwoTransPenaltyChangeLo{} to \STwoTransPenaltyChangeHi; Holm $p$~\STwoTransPenaltyChangePHolm), \STwoEssPenaltyChangeEst{} on essence change (\STwoEssPenaltyChangeLo{} to \STwoEssPenaltyChangeHi; Holm $p$~\STwoEssPenaltyChangePHolm) and \STwoCreatPenaltyChangeEst{} on creativity (\STwoCreatPenaltyChangeLo{} to \STwoCreatPenaltyChangeHi; Holm $p$~\STwoCreatPenaltyChangePHolm). Study~2's ordinal check (registered secondary), a model that treats the five answers as ordered categories rather than evenly spaced numbers, leaned toward dilution on transformativeness and essence change, not significantly (Holm $p$~\STwoTransPenaltyChangeOrdPHolm{} on both). On creativity its fit did not converge, so that check is not completed. With ratings at None near the floor, this test says little either way.

Study~4 registered the penalty at each level as secondary tests. At Slight, as in Study~5 (main text), they found a penalty that was significant without correction and not shown inside the bound on transformativeness (\SFourTransPenaltySlightEst{} points, $p$~\SFourTransPenaltySlightP) and essence change (\SFourEssPenaltySlightEst, $p$~\SFourEssPenaltySlightP), but not on creativity, where the Slight penalty, \SFourCreatPenaltySlightEst{} points, lay inside the bound. The essence-change difference at Slight had been reported before registration: from our February 2025 analysis of these data, the CogSci 2025 paper reported a label difference in essence change at Slight and at no other level.\cite{kamper2025originalsin}

\begin{figure}[htbp]
\centering
\makebox[\textwidth][c]{\includegraphics[width=183mm]{figS_pooled_penalty}}
\caption{\textbf{Exploratory: the gap on second works, pooled over studies.} A pooled statistical model fitted Studies 2, 4 and 5 together, each study weighted equally, and estimated the gap (the Human-label mean minus the AI-label mean) at each edit level for \textbf{a},~transformativeness, \textbf{b},~essence change and \textbf{c},~creativity of the second work. Filled squares pool all three studies; open squares leave out Study~4, whose creativity estimates lay below zero at every level, none significantly. Lines are 95\% confidence intervals, not corrected for multiple comparisons. Right of zero, the AI label lowered the rating; for essence change, the AI-label group judged the edit to change the original less. None is the original shown again unchanged, and ratings there sit near the floor of the scale. The fits converged but are singular: one source of variation that the model allows for was estimated at zero. The analysis was specified on 9 October 2026, after every registered result was known, and it changes no registered verdict.}
\label{figS:pooled}
\end{figure}

\paragraph*{The studies disagree.}
The pooled curve averages over studies that disagree on creativity (main-text Fig.~2b; Supplementary Table~\ref{tab:supp-labeleffects}). In each study's own model, Study~2's penalty grew from None to Dramatic (\STwoCreatPenaltyNoneEst, 95\% CI \STwoCreatPenaltyNoneLo{} to \STwoCreatPenaltyNoneHi; \STwoCreatPenaltySlightEst, \STwoCreatPenaltySlightLo{} to \STwoCreatPenaltySlightHi; and \STwoCreatPenaltyDramaticEst, \STwoCreatPenaltyDramaticLo{} to \STwoCreatPenaltyDramaticHi; exploratory and uncorrected, since Study~2 registered no test at each level). Study~4's lay below zero at every level (\SFourCreatPenaltyNoneEst, \SFourCreatPenaltySlightEst{} and \SFourCreatPenaltyDramaticEst; registered secondary, none significant), and Study~5's peaked at Slight (\SFiveCreatPenaltyNoneEst, \SFiveCreatPenaltySlightEst{} and \SFiveCreatPenaltyDramaticEst; registered secondary; at Dramatic, $p$~\SFiveCreatPenaltyDramaticP). In the pooled model, the creativity curvature was larger in Studies 4 and 5 (\XSCreatCurvatureSFourEst, uncorrected $p$~\XSCreatCurvatureSFourP, and \XSCreatCurvatureSFiveEst, $p$~\XSCreatCurvatureSFiveP) than in Study~2 (\XSCreatCurvatureSTwoEst). Leaving out one study at a time shows how much the pooled readings rest on each. Without Study~4, the pooled creativity penalty at None, Slight and Dramatic was \XSCreatPenaltyWithoutSFourNoneEst, \XSCreatPenaltyWithoutSFourSlightEst{} and \XSCreatPenaltyWithoutSFourDramaticEst. At Slight its 90\% CI (\XSCreatPenaltyWithoutSFourSlightLoNinety{} to \XSCreatPenaltyWithoutSFourSlightHiNinety) extended past the bound, where with all three studies it had stayed inside (\XSCreatPenaltySlightLoNinety{} to \XSCreatPenaltySlightHiNinety), and at Dramatic its 90\% CI reached \XSCreatPenaltyWithoutSFourDramaticHiNinety. The transformativeness reading at Slight is as fragile: without Study~2 or without Study~5, the upper end of its 90\% CI was \XSTransPenaltyWithoutSTwoSlightHiNinety{} or \XSTransPenaltyWithoutSFiveSlightHiNinety, past the bound. These fits were also singular but converged.

\paragraph*{Comparing the accounts as models.}
Four models were fitted to the same ratings, each freer than the one before: severance (no label effect), taint (one penalty per study, the same at every level), dilution (each study's penalty changing in a straight line across the levels) and a free shape (any penalty at each level, as in the pooled model). Two standard scores for comparing models, AIC (Akaike's information criterion) and BIC (the Bayesian information criterion), reward a model for fitting the ratings and charge it for each added quantity it estimates; BIC charges more. On every measure, AIC favored the free shape and BIC favored severance, the model with no penalty at all. Expressed as weights, each model's share of support among the four, AIC gave the free shape \XSCreatFreeShapeAICWeight{} on creativity, \XSTransFreeShapeAICWeight{} on transformativeness and \XSEssFreeShapeAICWeight{} on essence change, and BIC gave severance \XSCreatSeveranceBICWeight, \XSTransSeveranceBICWeight{} and \XSEssSeveranceBICWeight{} after rounding. On creativity, a likelihood-ratio test, which asks whether a model's added quantities improve its fit, found that giving each study one constant penalty fitted better than no penalty ($p$~\XSCreatTaintLRTP). Study~4's constant penalty ran below zero, however, so this does not show one penalty shared by the three studies. An ordinal check of the pooled creativity penalty did not converge and is not completed, so it adds no evidence either way.

\paragraph*{Explanations of the peak at Slight.}
The cause of the peak at Slight is unknown, and no argument in the paper rests on it. Five explanations remain open. A slight edit may leave Alex's share of the result unclear, so the label may matter more there. Of the ten Slight versions, \XSImageCheckSlightReframedN{} were also cropped or resized, so Slight is not a pure small edit. Ratings at None sit near the floor, which leaves little room for a penalty to show there. Under the Human label the original was Alex's own earlier work, which may earn the second work a bonus, and no condition with a neutral original was run to check this. The peak may also be chance. Nor is the residue confined to Slight: Study~2's creativity penalty (\STwoCreatPenaltyDramaticEst, uncorrected $p$~\STwoCreatPenaltyDramaticP; exploratory) and Study~5's transformativeness penalty (\SFiveTransPenaltyDramaticEst, $p$~\SFiveTransPenaltyDramaticP; registered secondary, uncorrected) were significant at Dramatic.
